\section*{Chapter \ref{ch:b1:factor-markets} --- Factor Markets}
% Numbers recomputed in code/bachelor-1/16-factor-markets/exercises/solutions.py (tested).

\begin{solution}{exo:b1:factor-markets:vmp-ex}
$\partial q/\partial L = 5L^{-0.5}$, so $\mathrm{VMP} = 20/\sqrt L$. Setting it equal to 5: $\sqrt L = 4$, $L = 16$.
\end{solution}

\begin{solution}{exo:b1:factor-markets:read}
A much lower share for Nigeria: with $86.6\%$ of workers self-employed, most
labour income is recorded as mixed income, which a naive share leaves out.
For France, with $12.1\%$ self-employed, the naive and adjusted shares are
closer.
\end{solution}

\begin{solution}{exo:b1:factor-markets:mono-ex}
$\mathrm{MCL} = 2 + 0.2L = 20 - 0.1L$: $L = 60$ and $w = 2 + 6 = 8$ (MRP 14). Competitive:
$2 + 0.1L = 20 - 0.1L$, $L = 90$, $w = 11$.
\end{solution}

\begin{solution}{exo:b1:factor-markets:floor12}
12 is above the competitive wage, so the firm hires until $\mathrm{MRP} = 12$:
$L = 80$, up from 60 (but below 90).
\end{solution}

\begin{solution}{exo:b1:factor-markets:best-floor}
The competitive wage, 11: employment 90, 30 more jobs than under \omterm{def:b1:factor-markets:monopsony}{monopsony}.
\end{solution}

\begin{solution}{exo:b1:factor-markets:markdown4}
$1/4 = 0.25$: the MRP is $25\%$ above the wage, and the wage is $1/1.25 = 80\%$ of the
MRP.
\end{solution}

\begin{solution}{exo:b1:factor-markets:derived-ex}
$\eta_L = -[(1 - 0.7) \times 1 + 0.7 \times 2] = -1.7$. With less elastic output demand the
scale effect shrinks and labour demand becomes less elastic: at $\varepsilon = -1$,
$-(0.3 + 0.7) = -1$.
\end{solution}

\begin{solution}{exo:b1:factor-markets:rental}
$10000 \times (0.04 + 0.10) + 200 = 1600$ dalens a year. Computers depreciate fast and
their prices fall fast, so both $\delta$ and $-\dot P_K/P_K$ are large: the rental price is
a large share of the purchase price.
\end{solution}

\begin{solution}{exo:b1:factor-markets:corr}
Strongly negative: the more workers are self-employed, the more labour
income a naive share leaves out, so naive shares fall as self-employment
rises.
\end{solution}

\begin{solution}{exo:b1:factor-markets:derive}
See \cref{prop:b1:factor-markets:markdown}: $\mathrm{MCL} = w(1 + 1/\varepsilon_L) = \mathrm{MRP}$. As $\varepsilon_L \to \infty$,
$1/\varepsilon_L \to 0$ and $w = \mathrm{MRP}$.
\end{solution}

\begin{solution}{exo:b1:factor-markets:above}
At 13, above the competitive 11, the firm hires until $\mathrm{MRP} = 13$: $L = 70$. That
is 10 more than under \omterm{def:b1:factor-markets:monopsony}{monopsony} (60) but 20 fewer than in the competitive
market (90), where a floor of 13 would also give 70: above the competitive
wage, the floor has the same effect in both, but the starting points
differ.
\end{solution}

\begin{solution}{exo:b1:factor-markets:separate}
A monopsonist faces a rising supply to the firm itself: a firm that raises
its wage alone should gain many workers only slowly, and a firm that cuts
it should lose few. Evidence: the response of a firm's recruits and quits
to its own wage (a small response means market power); wages that are
lower where there are fewer employers in the same occupation, holding the
industry's demand fixed; and employment that rises after a moderate wage
floor. Under competition, a firm paying less than others loses all its
workers, and the number of employers in a market does not matter.
\end{solution}

\begin{solution}{pb:b1:factor-markets:mill}
\textbf{1.} $5 + L/15 = 25 - L/15$: $L = 150$, $w = 15$.

\textbf{2.} $\frac{\dd L}{\dd w}\frac wL = 15 \times 15/150 = 1.5$.

\textbf{3.} $\tfrac12 \times 150 \times (15 - 5) = 750$.

\textbf{4.} $\tfrac12 \times 150 \times (25 - 15) = 750$.

\textbf{5.} 1,500.

\textbf{6.} $w(L)L = 5L + L^2/15$, so $\mathrm{MCL} = 5 + 2L/15$.

\textbf{7.} $5 + 2L/15 = 25 - L/15$: $L = 100$ and $w = 5 + 100/15 = 11.67$.

\textbf{8.} $\mathrm{MRP} = 25 - 100/15 = 18.33$, \omterm{def:b1:factor-markets:markdown}{markdown} $18.33/11.67 - 1 = 0.571$.

\textbf{9.} $\varepsilon_L = 15 \times 11.67/100 = 1.75$, and $1/1.75 = 0.571$.

\textbf{10.} $\int_0^{100}(25 - L/15)\,\dd L - 11.67 \times 100 = 2500 - 333.3 - 1166.7 = 1000$.

\textbf{11.} $\tfrac12 \times 100 \times (11.67 - 5) = 333.3$.

\textbf{12.} $1500 - 1000 - 333.3 = 166.7$, the triangle between 100 and 150 workers.

\textbf{13.} The floor is between 11.67 and 15: the mill hires everyone willing
to work at 13, $15 \times (13 - 5) = 120$.

\textbf{14.} $15 \times 10 = 150$.

\textbf{15.} Above 15 the mill hires until $\mathrm{MRP} = 17$: $15 \times (25 - 17) = 120$.

\textbf{16.} Employment exceeds 100 for floors above 11.67 and below
$\mathrm{MRP}(100) = 18.33$.

\textbf{17.} $15 \times (25 - 20) = 75$: fewer than under \omterm{def:b1:factor-markets:monopsony}{monopsony}.

\textbf{18.} 750 each, as in Part I: workers gain $750 - 333.3 = 416.7$, the mill loses
$1000 - 750 = 250$, and the \omterm{def:b1:partial-equilibrium-welfare:dwl}{deadweight loss} of 166.7 disappears.

\textbf{19.} Under \omterm{def:b1:factor-markets:monopsony}{monopsony} the floor removes the reason to hold
employment down: each extra worker no longer raises the wage of the
others. In a competitive market at 12, the firm already hires until
$\mathrm{MRP} = w$, and a floor of 15 moves it up its MRP curve to fewer workers.

\textbf{20.} Mill 1 pays $w = 5 + (L_1 + L_2)/15$, so its wage bill rises by
$\mathrm{MCL}_1 = 5 + (L_1 + L_2)/15 + L_1/15$.

\textbf{21.} $25 - L_1/15 = 5 + (2L_1 + L_2)/15$; with $L_1 = L_2$, $20 = 4L_1/15$: 75 at each
mill, 150 in all, at $w = 5 + 10 = 15$.

\textbf{22.} $\mathrm{MRP} = 25 - 5 = 20$; \omterm{def:b1:factor-markets:markdown}{markdown} $20/15 - 1 = 1/3$.

\textbf{23.} Mill 1 alone moving its employment changes the wage by
$\dd w/\dd L_1 = 1/15$, so $\varepsilon = 15 \times 15/75 = 3$, and $1/3$ is the \omterm{def:b1:factor-markets:markdown}{markdown}.

\textbf{24.} With two price-taking mills, labour demand is $30(25 - w)$, equal
to $15(w - 5)$ at $w = 18.33$, with 200 workers. The second mill raised the wage
from 11.67 to 15 and cut the \omterm{def:b1:factor-markets:markdown}{markdown} from 0.571 to 0.333, but the wage is
still below the competitive 18.33.

\textbf{25.} The employment-maximising floor in the one-mill town is 15
dalens an hour, the competitive wage; it raises employment from 100 to 150,
50 more workers.
\end{solution}
